%% =====================================================================
%%  B3-T-10-05 -- what B3 contributes to "The Withdrawal"
%%  -------------------------------------------------------------------
%%  LAYER A: APPEND-ONLY. Written once at the entry's write, then left
%%  alone. If the source has more to say later, add a bullet -- do not
%%  rewrite. Material found ONLY here goes in the `unique` slot with
%%  @unique set to yes, which makes it undeletable (rule R10).
%% =====================================================================
%% @kind: source
%% @id: B3-T-10-05
%% @book: B3
%% @topic: T-10-05
%% @locator: Law 16, pp. 121--127
%% @status: distilled
%% @unique: no
%% @integrated: yes
%% @updated: 2026-09-19

\dossier{B3}{T-10-05}{\bookshort{B3}}{Law 16, pp. 121--127}

\begin{distilled}
  \item Too much presence makes a person common: the constantly available are taken for granted and lose respect.
  \item The Deioces case: the judge withdrew completely, let the Medes taste life without him, and returned to a throne.
  \item Create a pattern of presence and absence; starve the other person of your presence to force respect.
  \item Extend the law of scarcity to skills: what is rare and hard to find instantly rises in value.
  \item Reversal: withdrawal only works after presence is established --- leave too early and you are simply forgotten.
\end{distilled}
